\section{Conclusion}\label{sec:conclusion}
Under exchangeability, the predictive of an exact TFM is the full conditional of a Gibbs sampler over unknown labels,
and a real TFM approximates it. \ours{} turns this into one procedure for candidate sets and label proportions. In
Study~2 a single chain with TabICL~v2 outperformed the specialised neural and linear learners we ran on new datasets.
Averaging over starts
improved on single chains, and averaging over TFMs improved on a TFM drawn at random; in a registered supplement the
multi-TFM average also beat untuned non-neural baselines. Whether a chain stalls depends on the TFM and on where the
chain starts. Open problems: a label-free diagnostic of stalling, a cheaper way to escape a wrong start than averaging over many chains (starting
chains from a trained learner was enough in our exploratory analysis), response-map bounds with usable constants, and
feature-dependent weak labels.
